% 相关工作章节
% 说明：
% 1. 本文件设计为被主文档 \input{} 进来使用。
% 2. 内容来源于 04_paper.md 的 Related Works 部分。
% 3. 这里只做 Markdown 到 LaTeX 的语法迁移，不改写正文内容。

\section{Related Works}

NeRF 及后续辐射场方法验证了照片级新视角合成的有效性，但体渲染采样开销较高，不适合作为逐帧在线更新的直接地图后端。3D Gaussian Splatting 通过显式高斯原语和可微光栅化显著提高了渲染效率，并被 MonoGS、GS-SLAM、SplaTAM、Photo-SLAM 等在线系统用于持续建图。然而，这些系统大多仍以 3DGS 体高斯为基础，新高斯通常依赖深度、点间距或协方差等规则初始化；在在线优化步数有限时，初始化误差会影响新插入区域的早期渲染质量和地图累积稳定性。

为提升高斯地图的几何一致性，2D Gaussian Splatting、GaussianSurfels 和 PGSR 等方法将高斯表示进一步表面化，通过二维 surfel、深度畸变约束和法线一致性约束改善表面贴合能力。这类工作说明，显式表面方向比 3D 椭球更适合承载深度和法线等几何先验。但现有表面化高斯方法主要面向离线或静态重建，尚未系统讨论在增量式系统中如何筛选新点、初始化完整 surfel 属性，并在少量在线优化内快速融入地图。

在大尺度户外场景中，LiDAR-Inertial-Camera 融合为增量式高斯建图提供了更稳定的位姿和几何观测。Gaussian-LIC 利用 LiDAR 点和视觉三角化点初始化 3DGS 地图，GS-LIVM 与 GS-LIVO 进一步考虑稀疏 LiDAR、局部窗口优化和资源约束，Gaussian-LIC2 则引入深度补全先验改善高斯初始化。这些工作与本文问题设定最接近，但它们仍主要采用 3DGS 后端，学习型先验也多用于补点或补深度，没有进一步预测尺度、方向、颜色和不透明度等完整属性。本文在 LIVO 前端提供位姿的条件下，将 2DGS surfel 作为地图表示，并通过前馈网络一次性预测新 surfel 的完整属性，从而改善增量插入阶段的初始质量